\section{Discussion}
\label{sec:discussion}

\subsection{The Legal Argument: Shared Hierarchy as Anti-Gerrymandering}

\paragraph{The structural constraint argument.}
NestSection's central legal contribution is the observation that a shared
bisection hierarchy is an anti-gerrymandering mechanism by \emph{construction},
not by detection.  Once the trunk is fixed by GeoSection at the spine level,
the congressional, senate, and house maps are all constrained to respect the
same initial geographic cuts.  A partisan actor who wishes to pack a community
into a safe congressional district must use the same geographic bisection that
the senate and house maps will also use.  If the trunk bisects that community,
the pack is constrained to one side of the trunk for all three chambers.

This is fundamentally different from post-hoc gerrymandering detection methods
(ensemble analysis, partisan symmetry tests), which evaluate a map after it is
drawn.  NestSection reduces the feasible plan space before any partisan
optimization occurs.

\paragraph{Degrees-of-freedom reduction.}
Formally, the number of distinct geographic configurations compatible with the
spine constraint is smaller than the number compatible with three independent
maps.  In a state with $g$ trunk regions, the congressional map is constrained
to respect the $g$ trunk boundaries; only within each trunk region does the
congressional mapmaker have freedom.  The partisan optimizer's effective search
space is reduced by a factor related to the trunk-to-state ratio of geographic
complexity.

This is the ``mutual constraint'' property: because the senate mapmaker is also
bound to the same trunk, the congressional mapmaker cannot gerrymander in a way
that the senate mapmaker could then undo.  The constraint is symmetric and
enforced at the algorithmic level, not at the political level.

\paragraph{Comparison to independent maps.}
With three independent maps, a unified partisan majority can achieve what
we call \emph{coherent gerrymandering}: packing or cracking the same
geographic community at all three levels simultaneously.  NestSection makes
coherent gerrymandering structurally harder: a pack at the congressional level
propagates to the senate and house levels, reducing the remaining degrees of
freedom for optimization at those levels.  Conversely, a crack at the
congressional level (splitting a community between two congressional districts)
automatically splits that community at the senate and house levels in the same
geographic location.

\paragraph{Relation to state constitutional nesting requirements.}
Many state constitutions and redistricting statutes require that districts be
``composed of contiguous, compact'' territory (e.g., NC Const.\ art.\ II, \S\S 3, 5)
or that county lines be respected.  Some statutes require state house districts
to be subsets of senate districts (e.g., Arizona, Indiana).  NestSection
operationalizes such requirements as a \emph{cross-chamber} computational
constraint rather than a within-chamber compactness check.

The NestSection compatibility score $\sigma$ provides a precise measure of
whether such a requirement is feasible: a state with $\sigma = 0$ can satisfy
exact nesting without any population imbalance penalty; a state with
$\sigma = 96.8\%$ (Texas) cannot satisfy any meaningful nesting without large
boundary exceptions.

\paragraph{Elections Clause scope.}
The Elections Clause (Art.\ I, \S4) grants states authority to regulate the
``Times, Places and Manner'' of congressional elections, subject to congressional
override.  A federal statute \emph{requiring} NestSection for congressional maps
could face Art.\ I, \S4 challenges from states arguing that the spine constraint
intrudes on their redistricting discretion (cf.\ \textit{Rucho v.\ Common Cause},
588 U.S.\ 684 (2019), which declined to review partisan gerrymandering of
congressional maps under federal law, leaving the field to state courts and
state law).  This paper does not advocate for such a federal mandate.

NestSection's most natural application is to \emph{state legislative} maps
--- the senate and house tiers --- where the governing law is state
constitutional provisions on compactness, contiguity, and integrity of
political subdivisions, not Art.\ I, \S4.  States that voluntarily adopt
NestSection for their state legislative chambers, drawing those maps to be
consistent with a pre-existing federal congressional map, face no federal
constitutional barrier.  The compatibility score $\sigma$ operates on the
relationship among all three chamber sizes, but the structural constraint
it enforces on the senate and house maps is a matter of state law.

The federal constitutional question arises only in the distinct scenario where
Congress were to mandate that state congressional maps conform to a
three-chamber spine constraint --- a proposal this paper explicitly does not
make.  The statutory design proposal in this section is addressed to
\emph{state legislatures and independent redistricting commissions}, which
have plenary authority over state legislative districts and, within Supremacy
Clause limits, broad discretion over the manner of drawing congressional
districts under Art.\ I, \S4.

\paragraph{Statutory design: the $g \geq k$ threshold.}
We propose a two-tier statutory structure:
\begin{enumerate}
  \item \textbf{Mandatory strict nesting} for states where $\gcd(C,S,H) \geq 5$:
        Oregon ($g=6$) and Alabama ($g=7$) currently qualify.  These states
        can achieve exact hierarchical nesting (Mode~1) with no population
        tolerance.  A statute requiring Mode~1 for these states would be
        straightforwardly achievable with existing redistricting algorithms.
  \item \textbf{Best-effort nesting} (Mode~3 with $\tau_\text{pop} \leq 0.02$)
        for states where $\gcd(C,S,H) \geq 2$: 22 states currently qualify.
        The statute would require the published boundary-zone fraction $\tau$ to
        be minimized and publicly reported, creating an empirically verifiable
        standard for cross-chamber consistency.
\end{enumerate}

\paragraph{Judicial enforceability.}
Judicial enforceability of a NestSection requirement raises two distinct
questions.  First, who has standing to challenge a non-nested map?  The most
natural plaintiffs are political subdivisions (counties, municipalities) whose
territories span chamber boundaries inconsistently~--- they can argue that the
inconsistent boundaries create administrative burdens that serve no compelling
interest.  Second, what is the standard of review?  If NestSection is framed
as a procedural rule (the map must be derived from a mathematical process),
courts reviewing it apply rational-basis scrutiny, not strict scrutiny~--- the
bar is low.  A statute stating ``state legislative districts shall be derived
by refining the congressional factorization tree'' has a rational basis in
administrative consistency and resistance to manipulation.  Courts striking
down such a requirement would need to find it facially arbitrary, which the
Bimodality Gap Theorem makes difficult: the compatible states are a
mathematically defined class, not a politically selected one.

\subsection{Limitations}

\paragraph{Incompatible states.}
For 18 states with $g = 1$ and congressional seat counts $C \geq 5$, NestSection
reduces to three independent GeoSection runs with a shared census graph.  The
algorithm provides no geographic anchor beyond the state boundary.  Statutory
mandates for cross-chamber consistency in these states would require either
(a)~changing the senate chamber size to be compatible with $C$ (a major
political change), or (b)~tolerating large boundary zones that may undermine
the legal defensibility of the nesting constraint.

\paragraph{Population balance.}
The spine constraint interacts with the statutory $\pm 0.5\%$ population
equality requirement for congressional districts.  In a strict-nesting state,
the trunk regions must have populations that are simultaneously divisible
into $C/g$, $S/g$, and $H/g$ equal-population sub-districts.  Population
heterogeneity within trunk regions may require slight adjustments to trunk
boundaries (Mode~2 or Mode~3) even in nominally compatible states.
The empirical validation in Section~\ref{sec:evaluation} is theoretical;
full empirical confirmation requires running the NestSection pipeline on
actual census data.

\paragraph{Single-seat states.}
Eleven states have $C = 1$, which trivially achieves $g = 1 = \min(C,S,H)$
and $\sigma = 0$.  However, a single-seat congressional delegation has no
internal geographic structure to share with the state legislature.  The
``strict compatibility'' of these states is a mathematical artifact of the
formula rather than a meaningful cross-chamber constraint.  In practice,
single-seat states require the congressional district to be the entire state,
which makes any cross-chamber nesting trivial (the congressional district
contains all senate and house districts by definition).

\paragraph{Nebraska's unicameral legislature.}
Nebraska is unique among US states in having a unicameral legislature (49
senators, no separate house).  The nest.rs implementation uses $H = S = 49$
for Nebraska; in practice, NestSection would apply only to the two-chamber
relationship (congressional vs.\ senate), and the house tier is vacuous.

\paragraph{Senate-first vs.\ congressional-first recursion.}
NestSection as defined uses the GCD trunk and then recurses independently
within each trunk region for all three chambers.  An alternative is to define
the spine from the largest chamber (house) and coarsen upward.  This is the
Mode~2 design from the plan: using the senate as the spine and placing
congressional seats at coarsened senate groups.  The Mode~2 approach is
mathematically equivalent to Mode~1 when divisibility holds, but differs in
the within-region recursion order when it does not.  We defer a full
comparison to future work.

\paragraph{Approximate compatibility: Arizona and Indiana as illustrative cases.}
Two near-compatible states illustrate a continuum that the strict bimodality
threshold ($\sigma = 0$ vs.\ $\sigma \geq 50\%$) obscures.  Arizona has
$C = 9$, $S = 30$, $H = 60$, giving $\gcd(9, 30, 60) = 3$ and $\sigma = 67\%$.
The spine trunk $[3]$ would create a three-way top-level geographic split
shared across all chambers~--- a natural east/west/central Arizona division
(roughly: Phoenix metro, Tucson corridor, and rural east).  Each side would
hold 3 congressional, 10 senate, and 20 house districts.  This is a
substantive, non-trivial compatible structure with real geographic content,
even at $\sigma = 67\%$ (incompatible by strict definition).  It illustrates
that the strict bimodality threshold hides a continuum of approximate
compatibility that future work could formalise.

Indiana has $C = 9$, $S = 50$, $H = 100$, giving $\gcd(9, 50, 100) = 1$ and
$\sigma = 89\%$.  This is the hardest incompatible case: 9 congressional, 50
senate, and 100 house districts share no common factor.  Any NestSection spine
for Indiana would require either (a)~rounding congressional districts to 10
(impossible by apportionment), or (b)~accepting that state legislative maps
cannot be nested in the congressional structure.  Indiana represents the class
of states where NestSection applies only to the state legislative chambers
internally~--- senate nests into house: $\gcd(50, 100) = 50$, perfect
compatibility~--- without congressional linkage.

\subsection{Future Work}

\paragraph{Empirical pipeline validation.}
The most important near-term extension is running NestSection on actual census
data for the case-study states (Oregon, North Carolina, Texas) and measuring:
(a) the actual nesting violation $V$ achieved on census geography;
(b) the edge-cut premium vs.\ independent GeoSection runs;
(c) the partisan outcome distribution across 100 GeoSection seeds.

\paragraph{Gerrymander resistance hypothesis.}
Section~\ref{sec:intro} conjectures that NestSection reduces partisan variance
by approximately 30--50\% in moderate-compatibility states.  Testing this
hypothesis requires the ensemble-diagnostics infrastructure from
\texttt{redist-analysis::ensemble\_diagnostics} (R-hat, ESS, Hamming
autocorrelation).  A formal test would compare the standard deviation of
Democratic seat share across 100 NestSection seeds vs.\ 100 independent
GeoSection seeds for NC and WI.

\paragraph{Reapportionment stability.}
When a state gains or loses a congressional seat after reapportionment,
the GCD $g$ may change, altering which nesting mode applies.  Montana gained
a seat in 2020 (going from $C=1$ to $C=2$), preserving $\sigma=0$ because
$g=2$ still equals $\min = 2$.  Texas will likely gain seats after 2030,
possibly achieving a GCD $> 1$ with its senate.  A longitudinal analysis of
how $\sigma$ has changed across census decades would quantify the stability
of nesting constraints over reapportionment cycles.

\paragraph{VRA interaction.}
The NestSection trunk constraint interacts with Voting Rights Act
Section~2 compliance.  If a minority-majority community is bisected by the
trunk, it will be split at all three legislative levels simultaneously.  This
could harm VRA compliance if the split reduces the minority's ability to elect
representatives of their choice.  The \texttt{redist analyze --types bloc-voting}
Callais pipeline provides the tools to measure this interaction.  A NestSection
implementation should include a VRA preflight check (analogous to the
\texttt{callais\_preflight} function in the existing pipeline) to detect cases
where the spine bisects a protected community.
